\chapter{Calprotectin (Stool)}

\section*{What Is This Test?}
The fecal calprotectin test measures the concentration of calprotectin in stool. Calprotectin is a calcium- and zinc-binding heterodimer protein (S100A8/S100A9) that constitutes approximately 60\% of the cytosolic protein of neutrophils. It is released from neutrophils during degranulation and activation. Fecal calprotectin directly reflects the degree of intestinal neutrophilic inflammation. It is elevated in inflammatory bowel disease (Crohn disease and ulcerative colitis) and in any condition causing intestinal inflammation (infectious enterocolitis, diverticulitis, NSAID enteropathy, microscopic colitis). It is normal or mildly elevated in irritable bowel syndrome (IBS), which is a disorder of gut-brain interaction without significant inflammation. Fecal calprotectin is the single most useful non-invasive test for distinguishing IBD from IBS and for monitoring disease activity in patients with established IBD. Its two main clinical uses are: (1) triage of patients with chronic diarrhea or abdominal pain to select those who need colonoscopy (to rule out IBD); and (2) monitoring of disease activity and response to therapy in established IBD, where it correlates well with endoscopic mucosal healing.

\section*{Alternative Names}
\begin{itemize}
  \item Fecal Calprotectin
  \item Stool Calprotectin
  \item Fecal Calprotectin Assay
  \item IBD Screen
\end{itemize}
\section*{Principle}
Fecal calprotectin is measured from a single random stool sample (1--5 g) by quantitative ELISA (enzyme-linked immunosorbent assay) or automated chemiluminescent immunoassay (CLIA) using monoclonal antibodies against calprotectin. The result is reported in micrograms of calprotectin per gram of stool (mcg/g). The most commonly used cutoff is 50 mcg/g: values below 50 mcg/g are normal and make active IBD unlikely; values above 50 mcg/g suggest intestinal inflammation and warrant further evaluation. Higher cutoffs (120--250 mcg/g) improve specificity for IBD. In patients with established IBD, a value above 250 mcg/g correlates with active endoscopic disease, and a value below 120 mcg/g suggests mucosal healing. Fecal calprotectin is stable in stool for up to 7 days at room temperature, and the protein resists degradation by intestinal and bacterial proteases. It is more sensitive and specific than fecal lactoferrin (another neutrophil marker) for IBD detection.

\section*{History}
Calprotectin was discovered in 1980 as a protein in the cytosol of human granulocytes and was initially called L1 protein or MRP-8/MRP-14. Its role as a fecal marker of intestinal inflammation was established in the 1990s. In 1992, Røseth and colleagues demonstrated that fecal calprotectin was elevated in patients with IBD and correlated with disease activity \cite{roseth1997}. Clinical trials established that fecal calprotectin could reliably distinguish IBD from IBS: a systematic review by van Rheenen and colleagues (2010) showed a pooled sensitivity of 93\% and specificity of 96\% for IBD in patients with chronic diarrhea, reducing the need for colonoscopy by two-thirds \cite{vanrheenen2010}. Subsequent studies confirmed that fecal calprotectin correlates with endoscopic activity and predicts mucosal healing during therapy \cite{dhaens2012}. Fecal calprotectin is now recommended by international guidelines (ECCO, ACG) as a first-line test in patients with suspected IBD and for monitoring disease activity in established IBD, reducing unnecessary colonoscopies.

\section*{Why This Test Is Ordered}
\begin{itemize}
  \item Differentiate IBD (Crohn disease, ulcerative colitis) from IBS in patients with chronic abdominal pain, diarrhea, or altered bowel habits, thereby avoiding unnecessary colonoscopy when the level is normal
  \item Estimate the probability of IBD in primary care before referral to a gastroenterologist
  \item Monitor disease activity in patients with established IBD and assess response to medical therapy (e.g., biologics, aminosalicylates)
  \item Predict clinical relapse in IBD: a rising fecal calprotectin precedes endoscopic and clinical relapse by 3--6 months \cite{dhaens2012}
  \item Evaluate chronic diarrhea or suspected intestinal inflammation in children
  \item Help select which patients with unexplained gastrointestinal symptoms require urgent versus elective endoscopy
\end{itemize}

\section*{Primary vs Secondary Test}
Fecal calprotectin is a primary triage test for suspected IBD and a primary monitoring test for established IBD. It is a secondary/supplementary test when compared with colonoscopy with biopsies, which remains the gold standard for the diagnosis of IBD.

\section*{How This Test Is Performed}
The patient collects a small stool sample (approximately 100--500 mg, about the size of a pea or hazelnut) in a clean collection container. The sample is transported to the laboratory at room temperature (stable for up to 7 days) or refrigerated. In the laboratory, the stool is weighed, an extraction buffer is added, and the calprotectin concentration is measured by ELISA or a quantitative lateral-flow immunoassay. Results are typically available within 1--3 days (or within hours with point-of-care devices). Ideally, the sample should be collected within 3 days of the test and refrigerated if not processed immediately.

\section*{What Preparation Is Needed}
No fasting is required. Patients should not use NSAIDs for 2 weeks before the test (NSAIDs cause intestinal inflammation and elevate fecal calprotectin). Stool should not be collected during an episode of acute infectious diarrhea (gastroenteritis, food poisoning, or after travel), because infection itself elevates fecal calprotectin and produces false-positive results. Collect the sample when symptoms are present (active diarrhea) but not during menstruation. A single sample is sufficient in most cases; a second sample may be obtained if results are borderline.

\section*{Contraindications and Precautions}
\begin{itemize}
  \item Active gastrointestinal bleeding or fresh blood in the stool can falsely elevate fecal calprotectin
  \item NSAIDs, aspirin, and proton pump inhibitors may produce mild elevations
  \item Children under 4 years have higher normal values than older children and adults
  \item Acute infectious gastroenteritis can transiently raise calprotectin; defer testing until resolution
  \item Sensitivity is reduced for isolated small bowel Crohn disease; a normal value does not fully exclude small bowel inflammation
  \item Results are not useful for determining disease location or extent
\end{itemize}
\section*{Risks}
No risks are associated with stool collection. The main risk is a false-negative or false-positive result leading to delayed diagnosis (missed IBD) or unnecessary colonoscopy.

\section*{What Happens After the Test}
Results are available within 1--3 days. A normal result (<50 mcg/g) effectively excludes IBD and avoids colonoscopy, whereas a borderline (50--120 mcg/g) or elevated (>120 mcg/g) result prompts endoscopic evaluation or repeat testing. In patients with established IBD, serial measurements guide therapy escalation and predict relapse; repeat fecal calprotectin is typically measured every 3--6 months during disease monitoring.

\section*{Understanding Results}
\subsection*{Normal (<50 mcg/g)}
Minimal or no neutrophilic intestinal inflammation. IBS is likely; IBD is unlikely (negative predictive value above 95\% for IBD). A normal fecal calprotectin level in a patient with chronic GI symptoms effectively excludes IBD and avoids unnecessary colonoscopy.
\subsection*{Borderline (50--120 mcg/g)}
Mild intestinal inflammation. May represent mild, treated, or quiescent IBD, NSAID enteropathy, microscopic colitis, infectious gastroenteritis, colorectal adenomas or carcinoma, or diverticulitis. Repeat testing and clinical correlation are indicated.
\subsection*{Elevated (120--250 mcg/g)}
Active intestinal inflammation. Consistent with IBD or other organic GI pathology. Colonoscopy with biopsies is typically warranted.
\subsection*{Markedly Elevated (>250 mcg/g)}
Significant neutrophilic intestinal inflammation. Highly suggestive of active IBD. Colonoscopy with biopsies is indicated. In a patient with known IBD, a value above 250 mcg/g indicates active mucosal inflammation and a high risk of clinical relapse within the next 3--6 months; escalation of IBD therapy should be considered.

\section*{Normal Results}
Normal: <50 mcg/g (negative). Borderline: 50--120 mcg/g. Elevated (IBD likely): 120--250 mcg/g. Markedly elevated: >250 mcg/g. Values are age-dependent: children under 4 years have higher normal fecal calprotectin levels. NSAIDs and aspirin can elevate fecal calprotectin (NSAID enteropathy). Proton pump inhibitors (PPIs) cause mild calprotectin elevation.

\section*{Follow-up Tests}
Colonoscopy with ileoscopy and biopsies (the gold standard for IBD diagnosis), upper endoscopy (for suspected Crohn disease with upper GI involvement), CT or MR enterography (for small bowel evaluation), fecal lactoferrin (an alternative fecal inflammatory marker), serum CRP and ESR, and repeat fecal calprotectin in 3--6 months (for IBD monitoring and to predict relapse).

\section*{References}
\bibliographystyle{unsrtnat}
\bibliography{references}

\clearpage
\section*{Regulatory Status and Authorization Date}
This chapter describes a test category, not a specific commercial product. FDA, CDSCO, EU IVDR, and MHRA approval or registration dates therefore cannot be assigned without the assay manufacturer, product name, instrument, and intended use. For laboratory-developed tests, an individual agency approval date may not exist. See the Preface for the interpretation of regulatory status.

\clearpage
